\documentclass[12pt]{article}
\usepackage[utf8]{inputenc}
\usepackage[T1]{fontenc}
\usepackage[margin=1in]{geometry}
\usepackage{setspace}
\usepackage{xcolor}
\usepackage[colorlinks=true, urlcolor=blue!60!black, linkcolor=blue!60!black]{hyperref}
\setlength{\parindent}{0em}
\setlength{\parskip}{0.6em}
\renewcommand{\baselinestretch}{1.0}
\begin{document}
\begin{center}
{\LARGE \textit{Craig v. Missouri}, 29 U.S. (4 Pet.) 410 (1830)\par}
\vspace{0.5em}
{\large Case Brief --- Two-Page Edition\par}
\end{center}

\section{Memory Jogger}
Missouri's interest-bearing loan-office certificates, issued on the state's faith to circulate as money, were unconstitutional ``bills of credit'' even though not legal tender (\href{https://supreme.justia.com/cases/federal/us/29/410/}{29 U.S. 410}).

\section{Detailed Case Facts}
After the Panic of 1819, Missouri authorized up to \$200,000 in small-denomination, interest-bearing certificates, lent to citizens, receivable for taxes, salaries, and salt, and backed by a redemption fund and ``the faith of the state.'' Hiram Craig and two others gave the State a \$199.99 note for certificates and defaulted.

\section{Procedural History}
Missouri sued on the note; the state courts ruled for Missouri; on writ of error the Supreme Court reversed.

\section{Judicial Votes}
4--3. Marshall wrote for Duvall, Story, and Baldwin; Johnson, Thompson, and McLean each dissented.

\section{Holding}
Paper issued by a state on its faith and intended to circulate as money is a bill of credit, whatever its label and even without tender status. A note given for such paper is void for illegal consideration.

\section{Analysis of Opinions}
Marshall found jurisdiction because the statute's validity was necessarily decided below. Drawing on colonial usage, he held that the clause targets state paper currency, not ordinary borrowing; that the Constitution prohibits things, not names; and that the separate tender clause shows bills of credit need not be legal tender. The dissenters stressed the interest, the pledged fund, and the absence of compelled circulation; McLean urged ``mutual forbearance.'' \textit{Briscoe} (1837) later upheld notes of a state-owned corporate bank.

\section{Examples: Future Applications}
\textbf{Same-Side:} small bearer certificates receivable for taxes and backed by state faith are invalid. \textbf{Opposite-Side:} large general-obligation bonds sold to investors are ordinary borrowing. \textbf{Fence-Sitter:} a fully reserved digital token issued by a state corporation.

\section{Critique}
\textit{Briscoe} let states evade \textit{Craig} through corporate form, testing its things-not-names principle; historians see Missouri's modest experiment as unlike reckless paper money, and the debtors received a windfall.

\section{Key Quotations}
\begin{quote}
``To emit bills of credit conveys to the mind the idea of issuing paper, intended to circulate through the community for its ordinary purposes as money, which paper is redeemable at a future day; this is the sense in which the terms have been always understood.''
\end{quote}
\noindent --- Chief Justice Marshall, Opinion of the Court, 29 U.S. at 432.

\end{document}
